\section{Introduction}
\label{sec:intro}
% Draft 2026-10-28, shortened from 2.4 to about 1.2 columns (long version
% below as comments). Framing: (1) causal localization pays off for the verdict
% target, (2) general assessment of generalization (Adrian's comment on the
% abstract). Motivation: cheap signals gate costly verification (literature
% checked 2026-10-28). Numbers: RQ answer boxes in experiments.tex,
% paired_shift.csv (fold=mean).
LLMs are increasingly used to judge code, including their own. They determine whether a candidate solves
its task, whether a change is safe to accept, and which of several candidates to
keep. Without external feedback, such judgments are unreliable. Models rarely
improve by verifying their own outputs~\cite{Huang2024SelfCorrect,
Kamoi2024SelfCorrection}, favor their own generations~\cite{Panickssery2024Self},
and remain inconsistent as code judges~\cite{Jiang2025CodeJudgeBench}. Reliable
verification is costly: executing tests requires an environment and a test
suite~\cite{Raghavendra2026Rubrics}, and reasoning about a candidate involves generating many additional
tokens. A cheap signal indicating when a judgment can be trusted could gate
these costlier checks rather than replace them~\cite{Rakasi2026Uncertainty,
jung2025trust}. Execution-free correctness predictors provide such a signal from
the model itself, such as the probability of its
answer~\cite{Spiess2025Calib, Campos2026Multicalib}, the agreement among sampled
solutions~\cite{Chen24Eigenscore, shi2026codemorethantext,
bouchard2026functional}, or its hidden states~\cite{bui2025openia,
Vu2026Autoprobe, ribeiro2026repe}. Probes on hidden states are particularly
cheap, since they reuse activations the model computes anyway, and are already
deployed as first-stage monitors~\cite{McKenzie2025HighStakes,
Kramar2026Probes}. Such a predictor serves two purposes: it can estimate whether
a program is correct, and whether the model's own verdict on it can be trusted.

Although predictors trained on hidden states are accurate in distribution, their
acknowledged weakness is distribution shift~\cite{Kramar2026Probes,
OrgadTGRSKB25}. For code this is critical, as a predictor fitted on one benchmark
or programming language will usually be applied to other projects and languages. For natural-language question answering, Huang et
al.~\cite{Huang2025Internal} proposed causal localization as a remedy: they use
counterfactual interventions, which copy part of the hidden state from one
input into another and observe how the model's answer changes, to find the
subspace that robustly determines the answer %, and report that predictors built on it
%are more robust 
under shift. Whether this promise holds for code correctness assessment 
%, and how the established predictors compare under shift, 
remains open. Instead, prior work on code has considered only a subset of predictor families
and shift settings, and has largely focused on predicting program
correctness~\cite{ribeiro2026robustness, he2026codecircuit, huang2025pttrust,
gros2025localized}. To our knowledge, none includes causal methods or asks whether the model's
own verdict can be trusted.

We therefore ask \emph{which execution-free signals of code correctness
survive a change of benchmark or programming language, and what explains those
that do}. We adapt two causal methods of Huang et al., \emph{counterfactual
simulation} and \emph{value probing}, to code and compare them with output
confidence, a probe on the hidden state, LAT~\cite{ribeiro2026repe} and
EigenScore~\cite{Chen24Eigenscore}. To do so, we use Qwen3-Coder-30B and gpt-oss-20b for candidate generation and as judges; we apply shifts between LiveCodeBench and BigCodeBench
and between Python, C++ and Java. We then answer four research questions: \textbf{RQ1} asks how well the methods tell
whether the model's verdict can be trusted, the target the causal methods were
designed for; \textbf{RQ2} asks which part of the causal methods carries their
signal; \textbf{RQ3} asks how well the methods predict whether the program is
correct, the target of most prior work; and \textbf{RQ4} asks whether
agreement across samples, which requires no supervised training, tells better than a single
pass which problems the model will solve. 

% Findings without numbers (2026-10-29): the numbers are in the abstract and
% RQ1/RQ3. Previous numeric version: 0.09 of 0.11 kept vs 0.03 (verdict,
% Qwen); +0.04 over the best non-causal method; program 0.05 of 0.12 (Qwen),
% 0.02 of 0.16 (gpt-oss).
Our results show that causal localization pays off where it was designed to:
under language shift, a probe on the localized subspace is the most robust
predictor of Qwen3-Coder's verdict, significantly ahead of every non-causal
method, and the learned subspace itself, not the intervention, carries this
gain. On gpt-oss, no method stands out. Beyond the causal methods,
generalization depends on the target and the type of shift: no verdict
predictor transfers reliably between benchmarks, program correctness
transfers largely because the predictors recognize difficult problems, and
agreement across samples gives no consistent advantage over a single pass.

% Under language shift, the methods differ in how much of their lead over the
% untrained output confidence they keep. On the verdict target, causal
% localization keeps most of it: for Qwen3-Coder, a probe on the localized
% subspace stays 0.09 AUROC ahead of confidence (0.11 in distribution), while a
% probe on the full hidden state keeps only 0.03. The subspace probe is significantly ahead of
% the best non-causal method ($+0.04$ AUROC on average) and ahead in each of the
% three data partitions. The learned subspace itself carries this gain, so a
% simple probe on it suffices; random subspaces of the same size do not provide
% it. On gpt-oss, the causal methods are level with the best non-causal method.
% On program correctness, the trained methods keep at most half of their lead
% (0.05 of 0.12 for Qwen3-Coder, 0.02 of 0.16 for gpt-oss). Between benchmarks,
% the picture reverses: no method transfers verdict prediction reliably, whereas
% program correctness transfers, largely because the methods recognize which
% problems the model fails on; within one problem every method is weak.
% Agreement across samples gives no consistent advantage over a single pass.


Our contributions are:
\begin{itemize}
  \item To our knowledge, the first comparison of execution-free correctness
  predictors for code across four families, including causal methods, under
  both benchmark and language shift and for two targets.
  \item An adaptation of causal localization~\cite{Huang2025Internal} to code
  correctness prediction.
  \item An analysis showing that the learned subspace carries the advantage of
  the causal methods, and a distinction between candidate-level prediction and
  problem-level reliability: pooled AUROC must be read alongside a predictor's
  ability to rank the candidates of one problem.
\end{itemize}

% ---------------------------------------------------------------------------
% Long version (2.4 columns, with Viola's edits; replaced 2026-10-28). Parts
% cut and where they live now: self-preference \cite{Panickssery2024Self}
% (dropped); "two purposes" (RQ1/RQ3); RQ display blocks (inline); probe loses
% three times as much (RQ1); protocol sentence (Study Design); DAS training
% measure does not indicate usefulness (RQ2); problem-only probe (RQ3);
% contribution-2 detail (Methods).
%
% LLMs increasingly judge code, including their own They determine whether a
% candidate solves its task, whether a change is safe to accept, and which of
% several candidates to keep. Without external feedback, such judgments are
% unreliable. Models rarely improve by verifying their own
% outputs~\cite{Huang2024SelfCorrect, Kamoi2024SelfCorrection}, favor their own
% generations~\cite{Panickssery2024Self}, and remain inconsistent as code
% judges~\cite{Jiang2025CodeJudgeBench}. Reliable verification is costly:
% executing tests requires an environment and a test
% suite~\cite{Raghavendra2026Rubrics}, and reasoning about a candidate involves
% generating many additional tokens. A cheap signal indicating when a judgment
% can be trusted could gate these costlier checks rather than replace
% them~\cite{Rakasi2026Uncertainty, jung2025trust}. Execution-free correctness
% predictors provide such a signal from the model itself, ... Such a predictor
% serves two purposes: it can estimate whether a program is correct, and
% whether the model's own verdict on it can be trusted.
%
% Although predictors trained on hidden states are accurate in distribution,
% their acknowledged weakness is distribution shift. For code this is
% critical: a predictor fitted on one benchmark or programming language will
% be applied to other projects and languages. For natural-language question
% answering, Huang et al. proposed causal localization as a remedy. They use
% counterfactual interventions to identify a subspace of the hidden state that
% causally determines the model's answer. According to the authors, a
% predictor based on this subspace is more robust under distribution shift.
% Whether this holds for code, and how the established predictors compare
% under shift, is open. Prior studies of execution-free correctness prediction
% for code examine one or two families of predictors, under one type of shift
% and for program correctness only. To our knowledge, no study has compared
% execution-free predictors across method families under both a change of
% benchmark and programming language. None includes causal methods or asks
% whether the model's own verdict can be trusted.
%
% We therefore ask: which execution-free signals of code correctness survive a
% change of benchmark or programming language, and what explains those that
% do? ... Four research questions follow from the overarching one. The causal
% methods model the judge's decision, so we begin with the target they were
% designed for.
% RQ1: How well can execution-free methods tell whether the model's own
%   verdict can be trusted, in distribution and under shift?
% RQ2: Which part of the causal methods carries their signal? Where they are
%   helpful, we separate the learned subspace from the intervention and
%   compare it with random subspaces.
% RQ3: How well do the methods predict whether the program is correct? This is
%   the target of most prior work related to code.
% RQ4: Does agreement across samples indicate better than a single pass which
%   problems the model will solve? Agreement requires no training and
%   therefore cannot degrade under shift.
%
% Findings: ... In contrast, a probe on the full hidden state loses three
% times as much. The advantage lies in the learned subspace itself: a simple
% probe on it is as robust as the full causal method, and thus as cheap as any
% probe, while random subspaces of the same size do not provide this
% improvement. ...
%
% Contributions: (1) To our knowledge, the first broad comparison ... No
% reported number depends on a layer chosen on the same data, the language
% experiments are repeated on three partitions of the problems, and methods are
% compared with paired significance tests. (2) Novel adaptation of causal
% localization to code correctness prediction. One subspace per layer is
% trained on candidates the model judges correctly, serving both targets. (3)
% Analysis: the learned subspace, not the intervention, carries the gain, so a
% simple probe on it suffices; tracing how well the subspace is learned does
% not indicate whether it is useful. (4) Two findings for the program
% correctness target: pooled AUROC largely reflects which problems the model
% solves, since a probe that has seen only the problem matches one that has
% seen the candidate; and under language shift, trained methods lose most of
% their lead over output confidence.

% ---------------------------------------------------------------------------
% Earlier notes (kept 2026-10-28; superseded by the draft above).
%
% Correctness estimation is a fundamental problem in machine learning
% evaluation and reliability assessment. It is the problem of predicting
% whether a model output matches the ground truth.
% In this work, we define code correctness as adherence to a given problem
% specification, which can be verified through test execution -- the standard
% measure used in code generation benchmarks and in practice.
% there are several approaches to correctness estimation of LLM generations,
% for natural language and code, with the best performing methods being based
% on internal representations
% But these methods are fragile under distribution shift
% Motivation: Steering agentic development (cheap signals for decisions within
% an agentic development loop)
% triage vs selection
% Comparable setup for introspective correctness prediction for code
% Nothing found shows probes degrading under benchmark or language shift for
% code. That gap is what our experiments fill, which is worth saying
% explicitly in the intro.
% Draft, overarching question (2026-10-24): We ask which execution-free
% signals of correctness survive benchmark and language shift, and what
% explains those that do. ...
% Übersichtsbild: Pipeline oder intervention
%
% From main.tex (2026-10-22 drafts):
% - Comparison of X execution-free correctness predictors from Y families, on
%   four (two?) models, five benchmarks in three programming languages, at
%   candidate level and at problem level.
% - Fragile under shift: internal representations have been used to estimate
%   the knowledge stored in LLMs, predict the risk of hallucination, predict
%   jailbreaking behaviours, and assess code correctness. However, a known
%   weakness of these approaches is their poor generalization behaviour across
%   tasks and datasets [OrgadTGRSKB25] -> focus on generalization.
% - 2 Questions: (1) Causal localization was proposed as a shift-robust
%   alternative for natural language [Huang2025Internal], is it beneficial for
%   code generation? (2) Across families, which degrades least, under which
%   type of shift (on which target)?
% - Draft contributions: comparison of six predictors from four families
%   (RQ1, RQ3, RQ4); adaptation of counterfactual simulation and value probing
%   to code, interventions at the last prompt token; analysis of where the
%   causal methods help (RQ2), incl. the withdrawn "the subspace encodes the
%   model's answer rather than the program's correctness".
